% =====================================================================
%  3. Arranjo institucional e visão geral da implementação
% =====================================================================
\section{Arranjo institucional e visão geral da implementação}
\label{sec:arranjo}

A ação será executada pelo Instituto Água e Terra (IAT), com coexecução do
Instituto de Desenvolvimento Rural do Paraná (IDR-Paraná) e da Companhia de
Saneamento do Paraná (Sanepar). O IAT responde pelo planejamento, pela
coordenação técnica, pelas contratações, pela fiscalização e pela validação
dos produtos. O IDR-Paraná apoia a mobilização social e o planejamento,
elabora os Planos Integrados de Propriedade (PIP) e atua diretamente na
frente de esgotamento sanitário das áreas de mananciais. A Sanepar executa,
com recursos próprios de contrapartida, o Programa Sanepar Rural. A
Secretaria das Cidades (SECID) participa do planejamento e da coordenação
técnica e valida a implantação dos modelos de gestão. Os municípios integram a
governança e fazem a articulação local, enquanto as comunidades beneficiadas
participam da gestão e da operação cotidiana dos sistemas \cite{mop2026}.

As contratações financiadas pelo Banco Mundial (BIRD) somam custo estimado de
US\$ 35,48 milhões, a serem executados em seis anos. A Sanepar participará
com contrapartida estimada em US\$ 12,31 milhões, destinada ao Programa
Sanepar Rural, que é executado por convênio com os municípios. As demais
atividades serão executadas por contratação pública, mediante licitação,
segundo as regras do Banco Mundial, ou diretamente pela equipe técnica do
IAT, no caso das atividades de planejamento, seleção e articulação
institucional \cite{mop2026}.

Toda contratação financiada pelo BIRD percorre o mesmo ciclo, representado na
\autoref{fig:ciclo-contratacao}. A equipe do IAT elabora o Termo de Referência
(TR) em cerca de três meses, até que ele esteja pronto para envio ao Banco
Mundial. Segue-se o processo licitatório, que se estende até o mês anterior
ao início da execução, quando o contrato é assinado e o serviço, a obra ou o
fornecimento começa. É esse ciclo que estrutura o cronograma de 2026 e 2027
apresentado no Anexo~\ref{an:cronograma}.

\begin{figure}[htbp]
  \centering
  \caption{Ciclo das contratações financiadas pelo BIRD}
  \label{fig:ciclo-contratacao}
  \fluxograma{ciclo-contratacao}
  \fonte{Elaborado pelo IAT (2026) com base no POA 2026--2027.}
\end{figure}

Na execução, as atividades se organizam em três blocos, que a
\autoref{fig:fluxograma-geral} apresenta de forma integrada. O primeiro bloco
é de estruturação institucional e social: a seleção dos municípios e
comunidades, a definição e a pactuação do modelo de gestão e o trabalho
técnico-social junto às famílias. O segundo bloco é a implantação física, que
se divide em duas frentes: o abastecimento de água, que inclui o Programa
Sanepar Rural, e o esgotamento sanitário, que inclui a frente do IDR-Paraná. O
terceiro bloco compreende a entrega dos sistemas, a capacitação dos
beneficiários, o monitoramento e o encerramento. Duas premissas orientam essa
sequência. A primeira é que as comunidades contempladas serão definidas por
meio do trabalho técnico-social. A segunda é que o modelo de gestão será
pactuado com as comunidades antes da instalação dos sistemas, de modo que
nenhuma obra seja entregue sem uma entidade preparada para operá-la
\cite{mop2026}. As seções seguintes descrevem, do começo ao fim, primeiro a
implementação do abastecimento de água e depois a do esgotamento sanitário.

\begin{figure}[htbp]
  \centering
  \caption{Visão geral da implementação do Subcomponente 3.2.b}
  \label{fig:fluxograma-geral}
  \fluxograma{fluxograma-geral}
  \fonte{Elaborado pelo IAT (2026) com base em \citeonline{mop2026}.}
\end{figure}

\pendencia{Os US\$ 35,48 milhões citados no MOP correspondem ao financiamento
BIRD de todo o Subcomponente~3.2, incluindo a política e o plano estadual
(3.2.a, R\$ 10,6 milhões). Pela planilha de custos, o financiamento do 3.2.b
é de R\$ 183,0 milhões, cerca de US\$ 33,5 milhões ao câmbio de R\$ 5,46
\cite{custo2026}. Corrigir o valor.}

\pendencia{O MOP atribui ao IDR-Paraná corresponsabilidade pelas
atividades de abastecimento (Quadro 30: ``IAT/IDR''), mas o texto não descreve
o papel do instituto nessa frente além da mobilização social. Detalhar.}
